% Transcription of Cayley's Collected Mathematical Papers, Volume 4
% PDF pages 251--275 (printed pages 229--253)
% Papers: 261 (conclusion, articles 41--57), 262 (articles 1--end of denumerate part)

\documentclass[11pt]{article}

% ============================================================
% PREAMBLE
% ============================================================
\usepackage[utf8]{inputenc}
\usepackage[T1]{fontenc}
\usepackage[english]{babel}
\usepackage{amsmath,amssymb}
\usepackage{geometry}
\usepackage{fancyhdr}
\usepackage{hyperref}
\usepackage{array}
\usepackage{longtable}
\usepackage[protrusion=true,expansion=false]{microtype}

\geometry{a4paper,margin=1in}

\pagestyle{fancy}
\fancyhf{}
\fancyhead[LE,RO]{\thepage}
\renewcommand{\headrulewidth}{0pt}

% Fallback providecommands
\providecommand{\mathbbm}[1]{\mathbb{#1}}
\providecommand{\sgn}{\operatorname{sgn}}

\setlength{\abovedisplayskip}{4pt plus 1pt minus 1pt}
\setlength{\belowdisplayskip}{4pt plus 1pt minus 1pt}

% Cayley delimiter convention shorthand
\newcommand{\Cdel}{\mid}

\hypersetup{
  pdftitle={Cayley Collected Papers Vol IV, pp.~229--253},
  pdfauthor={Arthur Cayley},
  hidelinks
}

\usepackage{graphicx}
\emergencystretch=3em
\tolerance=600
\begin{document}

\thispagestyle{empty}
\begin{center}
\vspace*{2cm}
{\Large\scshape The Collected Mathematical Papers}\\[0.4em]
{\large of}\\[0.4em]
{\Large\scshape Arthur Cayley}\\[1.5em]
{\large Volume IV --- pages 229--253}\\[0.6em]
{\normalsize Papers 261 (conclusion), 262}\\[2em]
\vfill
\end{center}
\newpage

\setcounter{page}{229}

% ============================================================
% PAPER 261 (conclusion) --- articles 41--57
% ON THE CONIC OF FIVE-POINTIC CONTACT AT ANY POINT OF A PLANE CURVE
% ============================================================

\noindent\textbf{261]}\hfill\textsc{contact at any point of a plane curve.}\hfill\textbf{229}

\medskip

\noindent 41. I have not sought to verify this theorem by my formul\ae. I remark, that combining it with the before-mentioned theorem, the five-pointic conic is completely determined as follows; viz.

\medskip

\noindent\textsc{Theorem.} The five-pointic conic touches the conic at the point of contact (two conditions); it passes through the two points in which the polar conic is intersected by the tangent to the cubic at the tangential of the point of contact (two conditions); and it passes through the point which is the third point of intersection with the cubic of the line joining the point of contact with its second tangential.

\medskip

\noindent 42. The construction for the point of simple intersection leads at once to that for the sextactic points; in fact, consider a point having for its tangential a point of inflexion: a point of inflexion is its own tangential, and the second tangential of the first-mentioned point is therefore the point of inflexion: the line joining the point with the second tangent is therefore the tangent at the point, and the point of simple intersection coincides with the point itself, that is, the point in question is a sextactic point.

\medskip

\noindent 43. I represent the equation of the five-pointic conic by
\[
  (a, b, c, f, g, h \Cdel X, Y, Z)^2 = 0;
\]
the value of $a$ is
\[
\begin{aligned}
  a ={}& 9(1 + 8l^3)\,x^4 y^2 z^2\\
       &{} + 3x^2 y^2 z^2\bigl(3l^2 y^2 - (1 + 2l^3) yz\bigr)(x^2 + 2lyz)\\
       &{} - Q(x^2 + 2lyz)^2,
\end{aligned}
\]
in which equation
\[
  Q = y^3 z^3 + z^3 x^3 + x^3 y^3 - l^3 x^2 y^2 z^2;
\]
or, reducing by the equation $x^3 + y^3 + z^3 + 6lxyz = 0$,
\[
  Q = y^3 z^3 + x^3(-x^3 - 6lxyz) - 3l^2 x^2 y^2 z^2,
\]
that is,
\[
  -Q = x^6 + 6lx^4 yz + 3l^2 x^2 y^2 z^2 - y^3 z^3,
\]
and we have
\[
\begin{aligned}
  a ={}& 9(1 + 8l^3)\,x^4 y^2 z^2\\
       &{} + 3x^2 y^2 z^2\bigl(3l^2 y^2 - (1 + 2l^3) yz\bigr)(x^2 + 2lyz)\\
       &{} + (x^6 + 6lx^4 yz + 3l^2 x^2 y^2 z^2 - y^3 z^3)(x^2 + 4xl^2 yz + 4l^2 y^2 z^2).
\end{aligned}
\]
We have in like manner
\[
\begin{aligned}
  2f ={}& 18(1 + 8l^3)\,b^2 y^2 z^2\\
        &{} + 3x^2 y^2 z^2\bigl[(3l^2 y^2 - (1 + 2l^3) zx)(z^2 + 2lxy) + (3l^2 z^2 - (1 + 2l^3) xy)(y^2 + 2lzx)\bigr]\\
        &{} - 2Q(y^2 + 2lzx)(z^2 + 2lxy);
\end{aligned}
\]
the coefficient of $3x^2 y^2 z^2$ in the second line is
\[
  = 6l y z^3 + (-1 + 4l^3) x(y^3 + z^3) - 4l(1 + 2l^3)\,x^2 yz;
\]
or reducing by the equation of the curve,
\[
  = (1 - 4l^3)\,x^4 + (2l - 32l^4)\,x^2 yz + 6l^2 x y^2 z^2;
\]

\newpage
\noindent\textbf{230}\hfill\textsc{on the conic of five-pointic}\hfill\textbf{[261}

\medskip

\noindent and the coefficient of $-2Q$ is
\[
  = y^2 z^2 + 2lx(y^3 + z^3) + 4l^2 x^2 yz;
\]
or reducing by the equation of the curve,
\[
  = -2lx^4 - 8l^2 x^2 yz + y^2 z^2.
\]
We thus have
\[
\begin{aligned}
  2f ={}& 18(1 + 8l^3)\,b^2 y^2 z^2\\
        &{} + 3x^2 y^2 z^2\bigl((1 - 4l^3)\,x^4 + (2l - 32l^4)\,x^2 yz + 6l^2 x y^2 z^2\bigr)\\
        &{} + 2(x^6 + 6lx^4 yz + 3l^2 x^2 y^2 z^2 - y^3 z^3)(-2lx^4 - 8l^2 x^2 yz + y^2 z^2).
\end{aligned}
\]
Reducing the expressions of $a$ and $2f$, we find for the coefficients $(a, b, c, f, g, h)$,
\[
\begin{aligned}
  a  ={}& x^{10} + 10lx^7 yz + 40l^2 x^4 y^2 z^2 + (5 + 120l^3)\,x^4 y^3 z^3 - 10lx^2 y^4 z^4 - y^5 z^5,\\[4pt]
  2f ={}& -4lx^{10} - 40l^2 x^7 yz + (5 - 120l^3)\,x^4 y^2 z^2 + 40lx^4 y^3 z^3 + 8l^2 x^2 y^4 z^4 - 2y^5 z^5,
\end{aligned}
\]
which gives the completely developed form of the equation of the five-pointic conic.

\medskip

\noindent 44. I investigate the coordinates of the point of simple intersection of the cubic and the five-pointic conic as follows: the equations of the two curves are
\[
\begin{aligned}
  X^3 + Y^3 + Z^3 + 6lXYZ &= 0,\\
  (a, b, c, f, g, h \Cdel X, Y, Z)^2 &= 0;
\end{aligned}
\]
or if we write
\[
\begin{array}{ll}
  \alpha = 1, & A = c,\\
  \beta = 6lXY, & B = 2(gX + fY),\\
  \gamma = X^3 + Y^3, & C = aX^2 + 2hXY + bY^2,
\end{array}
\]
then the two equations are
\[
\begin{aligned}
  \alpha Z^3 + \beta Z + \gamma &= 0,\\
  A Z^2 + B Z + C &= 0,
\end{aligned}
\]
and the result of the elimination of $Z$ will be
\[
  (\alpha Z_1^3 + \beta Z_1 + \gamma)(\alpha Z_2^3 + \beta Z_2 + \gamma) = 0,
\]
where $Z_1$, $Z_2$ are the roots of the equation $AZ^2 + BZ + C = 0$; that is, we have
\[
\begin{aligned}
  &\phantom{+{}} \alpha^2\,\tfrac{C}{A}\bigl(B^2 - 2AC\bigr)\\
  &{} + \alpha\beta\bigl(-BC + 2ABC\bigr)\\
  &\hphantom{xxxx}\,\alpha\gamma\\[2pt]
  &{} + \gamma^2\,CA^2\\
  &{} - \gamma\beta\,BA^2\\
  &{} + \beta^2\,A^2 = 0;
\end{aligned}
\]

\newpage
\noindent\textbf{261]}\hfill\textsc{contact at any point of a plane curve.}\hfill\textbf{231}

\medskip

\noindent and substituting for $A$, $B$, $C$ their values, but attending only to the terms which involve $X^7$ and $Y^7$, the result is
\[
  (a^3 + c^3 - 8f^3 + 6acf)\,X^7 + \ldots + (b^3 + c^3 - 8f^3 + 6bcf)\,Y^7 = 0.
\]

\medskip

\noindent 45. But the result of the elimination must obviously be
\[
  (Xy_1 - Yx_1)(Xy_2 - Yx_2) = 0,
\]
if $(x_1, y_1, z_1)$ are the coordinates of the point of simple intersection. Comparing the two results, and forming the analogous third equation, we may write
\[
\begin{aligned}
  x^7 x_1 &= b^3 + c^3 - 8f^3 + 6bcf,\\
  y^7 y_1 &= c^3 + a^3 - 8g^3 + 6cag,\\
  z^7 z_1 &= a^3 + b^3 - 8h^3 + 6abh,
\end{aligned}
\]
where the value of $x^7 x_1$ may also be written
\[
  (b + c - 2f)(b + c\omega - 2\omega^2 f)(b + c\omega^2 - 2\omega f),
\]
$\omega$ being an imaginary cube root of unity, and so for the other two terms. The factors of $x^7 x_1$ might be calculated from the identical equation
\[
\begin{aligned}
  bY^2 + 2f YZ + cZ^2 ={}& 9(1 + 8l^3)\,x^2 y^2 z^2\bigl(y^2 Y^2 + z^2 Z^2 + 2lx YZ\bigr)\\
  &{} + \bigl[(y^2 + 2lzx) Y + (z^2 + 2lxy) Z\bigr] x\\
  &{} \times \Bigl\{\,3x^2 y^2 z^2 \bigl[(3l^2 y^2 - (1 + 2l^3) zx) Y + (3l^2 z^2 - (1 + 2l^3) xy) Z\bigr]\\
  &\hphantom{{} \times \Bigl\{}\, - Q\bigl[(y^2 + 2lzx) Y + (z^2 + 2lxy) Z\bigr]\,\Bigr\}.
\end{aligned}
\]
I remark, that putting $x = 0$, we have
\[
  bY^2 + 2fYZ + cZ^2 = -y^2 z^2\,(yY + zZ)^2,
\]
and hence writing $1$ for $Y$, and $-1$, $-\omega$, $-\omega^2$ for $Z$, we have
\[
\begin{aligned}
  b + c - 2f          &= -y^2 z^2 (y - z)^2,\\
  b + c\omega^2 - 2\omega f   &= -y^2 z^2 (y - \omega z)^2,\\
  b + c\omega   - 2\omega^2 f &= -y^2 z^2 (y - \omega^2 z)^2;
\end{aligned}
\]
and hence the product of the three factors is
\[
  -y^6 z^6 (y - z)^2 (y - \omega z)^2 (y - \omega^2 z)^2,
\]
which is equal to $-y^6 z^6 (y^3 - z^3)^2$, or $-y^6 z^6 (y^3 + z^3)^2$, which vanishes in virtue of the assumed equation $x = 0$. This shows that the function $b^3 + c^3 - 8f^3 + 6bcf$ contains the factor $x$. I have not verified \emph{\`a posteriori}, but I assume it to be true, that it contains in fact the factor $x^7$, and consequently that the expressions for $x_1$, $y_1$, $z_1$ are rational and integral functions of $(x, y, z)$ of the degree $25$, and containing respectively the factors $x$, $y$, $z$.

\medskip

\noindent 46. In the theory of the cubic, a point which depends linearly upon a given point may be termed a derivative of such point. According to a very beautiful theorem of Professor Sylvester's, the coordinates of a derivative point are necessarily rational and integral functions of a square degree of the coordinates $(x, y, z)$ of the given point; and moreover, there is but one derivative point having its coordinates of any given square degree $m^2$, or, as we may express it, only one derivative point of the degree $m^2$. The successive tangentials are derivative points of the degrees $4$, $16$, $64$, \&c.; the third point of intersection with the cubic, of the line joining two

\newpage
\noindent\textbf{232}\hfill\textsc{on the conic of five-pointic}\hfill\textbf{[261}

\medskip

\noindent derivative points of the degrees $m^2$ and $n^2$ respectively, is a derivative point of the degree $(m + n)^2$. Thus the third point of intersection with the cubic, of the line joining the given point with its second tangential, is a derivative point of the degree $(4 + 1)^2$, and it is easy to see that the degree is not $9$; it is therefore $25$. The point of simple intersection of the five-pointic conic is a derivative point of the degree $25$; it is therefore, according to Professor Sylvester's general theory, identical with the point given by the former construction; this agrees with the before-mentioned theorem of Mr~Salmon.

\bigskip

\begin{center}
\textsc{IV. Independent investigation for the Cubic.}
\end{center}

\medskip

\noindent 47. The following is, in substance, the method by which I first obtained the equation of the five-pointic conic, for the cubic
\[
  X^3 + Y^3 + Z^3 + 6lXYZ = 0.
\]
Write for shortness
\[
\begin{aligned}
  U &= x^3 + y^3 + z^3 + 6lxyz,\\
  V &= (x^2 + 2lyz) X + (y^2 + 2lzx) Y + (z^2 + 2lxy) Z,\\
  W &= (X^2 + 2lYZ) x + (Y^2 + 2lXZ) y + (Z^2 + 2lXY) z,\\
  T &= X^3 + Y^3 + Z^3 + 6lXYZ,\\
  P &= ax + by + cz,\\
  \Pi &= aX + bY + cZ;
\end{aligned}
\]
then, $X$, $Y$, $Z$ being current coordinates, and $x$, $y$, $z$ the coordinates of a point of the cubic (so that $U = 0$), the equation of the cubic will be
\[
  T = 0,
\]
and the equation of a conic having with it an ordinary (two-pointic) contact at the point $(x, y, z)$, will be\footnote{I have introduced the factor 2, to make this correspond with the form $DW - \Pi DU = 0$, in the case in question, $m = 3$.}
\[
  2W - \Pi V = 0.
\]

\medskip

\noindent 48. Now imagine from the point of contact lines drawn to the other four intersections of the two curves; in the case of the five-pointic conic, three of these lines will coincide with the tangent $V = 0$, and the remaining line will be the line joining the point of contact with the point of simple intersection. The equations of the lines in question can be found by Joachimsthal's theorem, viz. if $(x, y, z)$ be the coordinates of a given point, and $(X, Y, Z)$ current coordinates, then if in the equations of any two curves we substitute for the coordinates, $\lambda x + \mu X$, $\lambda y + \mu Y$, $\lambda z + \mu Z$, and between the equations so obtained eliminate $\lambda$, $\mu$, the resulting equation will be that

\newpage
\noindent\textbf{261]}\hfill\textsc{contact at any point of a plane curve.}\hfill\textbf{233}

\medskip

\noindent of the lines drawn from the point $(x, y, z)$ to the points of intersection of the two curves. The point $(x, y, z)$ is any point whatever, and it may therefore be a point of intersection, or, as in the present instance, a point of contact of the two curves; the only difference is, that in either case the degree of each equation as regards $(\lambda, \mu)$ is reduced by unity, and the degree of the resulting equation in $X$, $Y$, $Z$ is also reduced by unity: in the case of a point of simple intersection this is the only reduction; but in the case of a point of contact, the resulting equation contains the equation of the tangent as a factor, and rejecting this factor, the reduction in degree is two units.

\medskip

\noindent 49. Applying the method to the two equations, $T = 0$, $2W - \Pi V = 0$, and substituting therein for the original current coordinates $X$, $Y$, $Z$ the values $\lambda x + \mu X$, $\lambda y + \mu Y$, $\lambda z + \mu Z$, the equations become
\[
\begin{aligned}
  \lambda^3 U + 3\lambda^2 \mu V + 3\lambda \mu^2 W + \mu^3 T &= 0,\\
  2(\lambda^2 U + 2\lambda \mu V + \mu^2 W) - (\lambda P + \mu \Pi)(\lambda V + \mu V) &= 0;
\end{aligned}
\]
or writing $U = 0$, and omitting from each equation the factor $\mu$, the equations become
\[
\begin{aligned}
  3\lambda^2 V + 3\lambda \mu W + \mu^2 T &= 0,\\
  \lambda(4 - P) V + \mu(2W - \Pi V) &= 0;
\end{aligned}
\]
and putting in the first equation $\lambda = 2W - \Pi V$, $\mu = -(4 - P) V$, the result of the elimination contains the factor $V$, rejecting which it becomes
\[
  3(2W - \Pi V)^2 - 3(4 - P)(2W - \Pi V) W + (4 - P)^2 V T = 0,
\]
which is of the fourth degree in $(X, Y, Z)$, as it should be, and represents therefore the lines drawn from the point of contact to the other four points of intersection of the conic and cubic.

\medskip

\noindent 50. The equation may be written
\[
  -3(2W - \Pi V)\bigl((2 - P) W + \Pi V\bigr) + (4 - P)^2 V T = 0,
\]
and we obtain at once the condition that this may contain the factor $V$, viz. this condition is
\[
  P = 2;
\]
and if this be satisfied the conic will have a three-pointic contact, and there will be three other points of intersection. And writing $P = 2$, and throwing out the factor $V$, we find
\[
  3\Pi^2 V - 6\Pi W + 4T = 0
\]
for the equation of the lines from the point of contact to the three points of intersection. And we have now to determine $\Pi$ so that the function on the left hand may divide by $V^2$.

\medskip

\hfill {\footnotesize\textsc{c. iv.}\hfill 30}

\newpage
\noindent\textbf{234}\hfill\textsc{on the conic of five-pointic}\hfill\textbf{[261}

\medskip

\noindent 51. I simplify my original method by the use of a theorem of Mr~Salmon's, viz. writing
\[
\begin{aligned}
  T &= X^3 + Y^3 + Z^3 + 6lXYZ,        & \mathfrak{H} &= l^2(X^3 + Y^3 + Z^3) - (1 + 2l^3) XYZ,\\
  W &= (X^2 + 2lYZ) x + \ldots,         & \mathfrak{H}_1 &= \bigl(3l^2 X^2 - (1 + 2l^3) YZ\bigr) x + \ldots,\\
                                       & & H_1 &= \bigl(3l^2 x^2 - (1 + 2l^3) yz\bigr) X + \ldots,\\
  U &= x^3 + y^3 + z^3 + 6lxyz,         & H &= l^2(x^3 + y^3 + z^3) - (1 + 2l^3) xyz,
\end{aligned}
\]
we have identically
\[
  TH - U\mathfrak{H} = W H_1 - V \mathfrak{H}_1;
\]
and in the present case, since $U = 0$,
\[
  TH = W H_1 - V \mathfrak{H}_1.
\]
Hence, multiplying by $H$ and substituting this value of $TH$, the equation becomes
\[
  3H \Pi^2 V - 6H \Pi W + 4 W H_1 - 4 V \mathfrak{H}_1 = 0,
\]
or, as we may write it,
\[
  V\bigl(3H \Pi^2 - 4\mathfrak{H}_1\bigr) + 2W\bigl(2H_1 - 3H \Pi\bigr) = 0;
\]
and we can at once make the equation divide by $V$, viz. by assuming
\[
  \Pi = \tfrac{2}{3}\,\tfrac{H_1}{H} + AV,
\]
where $A$ is arbitrary: we have thus a four-pointic contact. And substituting for $\Pi$, and throwing out the factor $V$, the equation becomes
\[
  3H\bigl(\tfrac{2}{3}\,\tfrac{H_1}{H} + AV\bigr)^2 - 4\mathfrak{H}_1 + 2W\bigl(\tfrac{2}{3}\,A\bigr) = 0;
\]
or reducing,
\[
  16(H_1^2 - 3H \mathfrak{H}_1) - 36 H^2 A W + 24 H H_1 A V + 9 H^2 A^2 V^2 = 0,
\]
which is the equation of the lines drawn from the point of contact to the remaining two points of intersection.

\medskip

\noindent 52. I write for greater convenience
\[
  \Theta = \tfrac{16}{9}\,\mathfrak{H},
\]
($\Theta$ being as yet indeterminate); the equation is thus reduced to
\[
  9 H^2 \bigl\{ H(H_1^2 - 3H \mathfrak{H}_1) + \Theta W \bigr\} - 6 H^2 H_1 \Theta V + \Theta^2 V^2 = 0,
\]
and we have then to determine $\Theta$ so that the left-hand side may divide by $V$; or, what is the same thing, we must determine $\Theta$ so that
\[
  H(H_1^2 - 3H \mathfrak{H}_1) + \Theta W
\]

\newpage
\noindent\textbf{261]}\hfill\textsc{contact at any point of a plane curve.}\hfill\textbf{235}

\medskip

\noindent may divide by $V$. This implies the existence of an identical equation,
\[
  H(H_1^2 - 3H \mathfrak{H}_1) + \Theta W = M U + N V,
\]
which for $U = 0$ would give the decomposition in question; but I have not investigated the values of $M$ and $N$. I assume at the outset $U = 0$, and putting as before
\[
  -Q = x^6 + 6lx^4 yz + 3l^2 x^2 y^2 z^2 - y^3 z^3,
\]
and writing also
\[
\begin{aligned}
  S ={}& Xx\bigl[(1 + 8l^3)\,x^4 + (4l + 4ll^4)\,x^2 yz + (-2l^2 + 2l^5)\,y^2 z^2\bigr]\\
       &{} + Yy\bigl[(1 + 8l^3)\,y^4 + (4l + 4ll^4)\,y^2 zx + (-2l^2 + 2l^5)\,z^2 x^2\bigr]\\
       &{} + Zz\bigl[(1 + 8l^3)\,z^4 + (4l + 4ll^4)\,z^2 xy + (-2l^2 + 2l^5)\,x^2 y^2\bigr],
\end{aligned}
\]
I remark that for $U = 0$ we have
\[
  -xyz\,H_1^2 + V S + (1 + 8l^3)(QW - 3x^2 y^2 z^2\,\mathfrak{H}_1) = 0,
\]
an equation which, observing that
\[
  H = -(1 + 8l^3)\,xyz,
\]
and assuming also
\[
  \Theta = (1 + 8l^3)^2 Q,
\]
may be written
\[
  H(H_1^2 - 3H \mathfrak{H}_1) + \Theta W = -(1 + 8l^3) V S,
\]
which gives the required decomposition, viz., $\Theta$ having this value, the conic will have a five-pointic contact. Reducing by the last equation, and throwing out the factor $V$, we find
\[
  -9H^2(1 + 8l^3)\,S - 6 H^2 H_1 \Theta + \Theta^2 V = 0
\]
for the equation of the line joining the point of contact with the point of simple intersection. And if in this equation we write $H = -(1 + 8l^3)\,xyz$, and $\Theta = (1 + 8l^3)^2 Q$, we obtain, finally, for the equation of the line in question,
\[
  9x^2 y^2 z^2 S - 6 x^2 y^2 z^2\,Q H_1 + Q^2 V = 0,
\]
which is the before-mentioned result.

\medskip

\noindent 53. It only remains to verify the assumed equation
\[
  -xyz\,H_1^2 + V S + (1 + 8l^3)(Q W - 3 x^2 y^2 z^2\,\mathfrak{H}_1) = 0.
\]
We may write
\[
  -Q = x^6 + 6lx^4 yz + 3l^2 x^2 y^2 z^2 - y^3 z^3,
\]
and then observing that
\[
\begin{aligned}
  W &= xX^2 + \ldots + 2lx YZ + \ldots,\\
  \mathfrak{H}_1 &= 3l^2 x X^2 + \ldots - (1 + 2l^3) x YZ - \ldots,
\end{aligned}
\]

\medskip

\hfill {\footnotesize 30--2}

\newpage
\noindent\textbf{236}\hfill\textsc{on the conic of five-pointic}\hfill\textbf{[261}

\medskip

\noindent we find at once
\[
\begin{aligned}
  (1 + 8l^3)\bigl(QW - 3x^2 y^2 z^2\,\mathfrak{H}_1\bigr) ={}& -(1 + 8l^3)\Bigl\{ (x^7 + 6lx^5 yz + \ldots) X^2\\
  &\hspace{4em}{} + (2lx^7 + 12l^2 x^5 yz - 3x^2 y^3 z^2 - 2lx^4 y^2 z^2)\,YZ + \ldots \Bigr\}.
\end{aligned}
\]
Next writing
\[
\begin{aligned}
  H_1 &= \bigl(3l^2 x^2 - (1 + 2l^3) yz\bigr) X + \ldots,\\
  V   &= (x^2 + 2lyz) X + \ldots,\\
  S   &= x\bigl\{ (1 + 8l^3)\,x^4 + (4l + 4ll^4)\,x^2 yz + (-2l^2 + 2l^5)\,y^2 z^2 \bigr\} X + \ldots,
\end{aligned}
\]
and forming the expression for
\[
  -xyz\,H_1^2 + V S,
\]
the coefficient of $X^2$ is
\[
\begin{aligned}
  & -xyz\bigl[3l^2 x^2 - (1 + 2l^3) yz\bigr]^2\\
  &\quad + (x^2 + 2lyz)\,x\bigl\{ (1 + 8l^3)\,x^4 + (4l + 4ll^4)\,x^2 yz + (-2l^2 + 2l^5)\,y^2 z^2 \bigr\},
\end{aligned}
\]
which is
\[
  = (1 + 8l^3)\bigl(x^7 + 6lx^5 yz + 12l^2 x^3 y^2 z^2 - x y^3 z^3\bigr);
\]
the coefficient of $YZ$ is
\[
\begin{aligned}
  & -2xyz\bigl(3l^2 y^2 - (1 + 2l^3) zx\bigr)\bigl(3l^2 z^2 - (1 + 2l^3) xy\bigr)\\
  &\quad + (y^2 + 2lzx)\,z\bigl[(1 + 8l^3)\,z^4 + (4l + 4ll^4)\,z^2 xy + (-2l^2 + 2l^5)\,x^2 y^2\bigr]\\
  &\quad + (z^2 + 2lxy)\,y\bigl\{ (1 + 8l^3)\,y^4 + (4l + 4ll^4)\,y^2 zx + (-2l^2 + 2l^5)\,x^2 y^2 \bigr\},
\end{aligned}
\]
which is
\[
  = (1 + 8l^3)\bigl\{ y^2 z^2(y^3 + z^3) + 8lxy^3 z^3 + 12l^2 x^2 yz\,(y^3 + z^3) - 2lx(y^3 + z^3) - 2xy^2 z^2 \bigr\};
\]
or substituting for $y^3 + z^3$ and $y^3 + z^3$ the values $-x^3 - 6lxyz$ and $(x^3 + 6lxyz)^2 - 2y^3 z^3$ respectively, this is
\[
  = (1 + 8l^3)\bigl(2lx^7 + 12l^2 x^5 yz - 3 x^2 y^3 z^2 - 2lx^4 y^2 z^2\bigr).
\]
We have thus
\[
  -xyz\,H_1^2 + V S = (1 + 8l^3)\Bigl\{ (x^7 + 6lx^5 yz + 12l^2 x^3 y^2 z^2 - x y^3 z^3) X^2 + (2lx^7 + 12l^2 x^5 yz - 3 x^2 y^3 z^2 - 2lx^4 y^2 z^2) YZ + \ldots \Bigr\}
\]
and the equation
\[
  -xyz\,H_1^2 + V S + (1 + 8l^3)(Q W - 3 x^2 y^2 z^2\,\mathfrak{H}_1) = 0
\]
is thus verified.

\medskip

\noindent\textit{Addition.}---The foregoing memoir was communicated to Mr~Salmon, and I am indebted to him for two notes, containing the extension to a curve of any order, of the preceding investigation for the case of a cubic; I reproduce this extension in the following section.

\newpage
\noindent\textbf{261]}\hfill\textsc{contact at any point of a plane curve.}\hfill\textbf{237}

\bigskip

\begin{center}
\textsc{V. Extension of the last preceding method to a curve of any order.}
\end{center}

\medskip

\noindent 54. Consider the curve of the $m$-th order $T = 0$, and in the place of the coordinates write $\lambda x + \mu X$, $\lambda y + \mu Y$, $\lambda z + \mu Z$, where, as before, $(x, y, z)$ are the coordinates of the point of the curve, and $X$, $Y$, $Z$ current coordinates; the term involving $\lambda^m$ vanishes, and dividing out the factor $\mu$ the equation becomes
\[
  \lambda^{m-1} DU + \tfrac{m-1}{1}\,\lambda^{m-2}\mu\,D^2 U + \tfrac{m-1\cdot m-2}{1\cdot 2}\,\lambda^{m-3}\mu^2\,D^3 U + \tfrac{1}{4}\lambda^{m-4}\mu^3\,D^4 U + \&\text{c.} = 0.
\]
Making the like substitution in $D^2 U - \Pi DU = 0$, the assumed equation of the five-pointic conic, the factor $\mu$ divides out and the equation becomes
\[
  \lambda(2 DU - P\,DU) + \mu(D^2 U - \Pi DU) = 0,
\]
or, what is the same thing,
\[
  \lambda\,DU\bigl(2(m - 1) - P\bigr) + \mu(D^2 U - \Pi DU) = 0;
\]
and if from the two equations we eliminate $\lambda$, $\mu$, the result, throwing out the factor $DU$, is
\[
  (D^2 U - \Pi DU)^{m-1} - \tfrac{1}{4}\bigl(2(m - 1) - P\bigr)\,D U\,(D^2 U - \Pi DU)^{m-2} + \&\text{c.} = 0,
\]
where all the terms after the second contain the factor $DU$; the condition in order that the equation may divide by $DU$, is consequently $2(m - 1) - P = 2$, or $P = 2(m - 2)$, the condition of a three-pointic contact. Substituting this value, and dividing by $DU$, the equation becomes
\[
  -\Pi(D^2 U - \Pi DU)^{m-2} + \tfrac{1}{3} D^3 U(D^2 U - \Pi DU)^{m-3} - \tfrac{1}{4} D^2 U\cdot DU(D^2 U - \Pi DU)^{m-4} + \&\text{c.} = 0,
\]
which will be divisible by $DU$ if $-\Pi D^2 U + \tfrac{1}{3} D^3 U$ is divisible by $DU$, and the condition for this is found to be, as before,
\[
  \Pi = \tfrac{1}{3}\,\tfrac{DH}{H} + A\,DU,
\]
where $A$ is arbitrary; we have thus the conditions of a four-pointic contact.

\medskip

\noindent 55. Substituting this value of $\Pi$, we see that $D^3 U - \tfrac{1}{H} DH\cdot D^2 U$ divides by $DU$, viz. there exists an identical equation,
\[
  D^3 U - \tfrac{1}{H} DH\cdot D^2 U = I U + J\,DU;
\]
and hence if $U = 0$,
\[
  \tfrac{1}{DU}\bigl(D^3 U - \tfrac{1}{H} DH\cdot D^2 U\bigr) = J,
\]
where $J$ is a quadric function of $(X, Y, Z)$. I do not know the general form of this function, but Mr~Salmon has obtained a result which may be generalised as

\newpage
\noindent\textbf{238}\hfill\textsc{on the conic of five-pointic}\hfill\textbf{[261}

\medskip

\noindent follows, viz. writing for $X$, $Y$, $Z$ the values $B\nu - C\mu$, $C\lambda - A\nu$, $A\mu - B\lambda$ (where, as before, $A$, $B$, $C$ are the first derived functions of $U$ and $\lambda$, $\mu$, $\nu$ are arbitrary), the expression for $J$ is
\[
  J = \tfrac{1}{(m-1)^2}\Bigl( D^2 U - \tfrac{1}{H} DH\cdot D^2 U \Bigr) = \tfrac{2(m - 2)}{(m - 1)}\,\Omega_3 - \tfrac{(m - 4) H}{(m - 1) H}\,\Omega_2,
\]
a formula which will be presently useful.

\medskip

\noindent 56. The foregoing equation may be written
\[
\begin{aligned}
  &(D^2 U)^{m-2}\bigl(-\Pi D^2 U + \tfrac{1}{3} D^3 U\bigr)\\
  &\quad + (D^2 U)^{m-3} DU\bigl\{(m - 2)\,\Pi^2 D^2 U - \tfrac{1}{3}(m - 3)\,\Pi D^3 U - \tfrac{1}{4} D^4 U\bigr\} + \&\text{c.}\,(DU)^2 \ldots = 0;
\end{aligned}
\]
and the term $-\Pi D^2 U + \tfrac{1}{3} D^3 U$ is equal to
\[
  \tfrac{1}{3}\bigl(D^3 U - \tfrac{1}{H} DH\cdot D^2 U\bigr) - A\,DU\cdot D^2 U, \qquad \tfrac{1}{3} J\,DU - A\,DU\cdot D^2 U.
\]
Substituting this value the equation divides by $DU$, and throwing out this factor it becomes
\[
\begin{aligned}
  &(D^2 U)^{m-2}\bigl(\tfrac{1}{3} J - A\,D^2 U\bigr)\\
  &\quad + (D^2 U)^{m-3}\bigl\{(m - 2)\,\Pi^2 DU - \tfrac{1}{3}(m - 3)\,\Pi D^3 U - \tfrac{1}{4} D^4 U\bigr\} + \&\text{c.}\,DU = 0;
\end{aligned}
\]
or observing that $\tfrac{1}{3}\Pi D^3 U = \Pi^2 D^2 U + \text{term containing } DU$, this may be written
\[
  (D^2 U)^{m-2}\bigl(\tfrac{1}{3} J - A\,D^2 U\bigr) + (D^2 U)^{m-3}\bigl(\Pi^2 D^2 U - \tfrac{1}{4} D^4 U\bigr) + \&\text{c.}\,DU = 0.
\]

\medskip

\noindent 57. If the equation divides by $DU$ we shall have a five-pointic contact; the condition for this is that
\[
  -A(D^2 U)^2 + \Pi^2 D^2 U + \Pi^2 D^2 U - \tfrac{1}{4} D^4 U
\]
may divide by $DU$, or more simply that
\[
  -A(D^2 U)^2 + \tfrac{1}{3} J\,D^2 U + \tfrac{1}{9}\bigl(\tfrac{DH}{H}\bigr)^2 D^2 U - \tfrac{1}{4} D^4 U
\]
may divide by $DU$, or, what is the same thing, the function in question must vanish in virtue of the substitution of the values $B\nu - C\mu$, $C\lambda - A\nu$, $A\mu - B\lambda$ in the place of $X$, $Y$, $Z$. The expression for $J$ has just been given; we have besides
\[
  \bigl(\tfrac{DH}{H}\bigr)^2 = \tfrac{m^2}{(m-1)^2}\,\Omega_1, \qquad D^2 U = \mu^2 \Omega_2, \qquad D^4 U = -\mu^4 \Omega_4,
\]
where the values of $\Omega_1$, $\Omega_2$ are given (ante, No.~27); we have thus

\newpage
\noindent\textbf{261]}\hfill\textsc{contact at any point of a plane curve.}\hfill\textbf{239}

\medskip
and
\[
  D^2 U = \tfrac{1}{(m-1)^2}\,\mu^2 \Omega_2,
\]
whence
\[
\begin{aligned}
  \Lambda ={}& \tfrac{1}{\mu^2}\Bigl[\tfrac{3(m-3)}{m-1}\,\Omega_4 - \tfrac{1}{(m-1)^2}\,\Omega_2^2\\
            &\qquad - \tfrac{1}{(m-1)^2}\,\Omega_2 \mu^2\Bigl\{\tfrac{4(m-2)}{m-1}\,\Phi + \tfrac{(m-2)}{(m-1) H}\,\Omega_2 - \tfrac{1}{(m-1)^2}\,\tfrac{1}{H}\,\Phi\,\mu^2\Bigr\}\\
            &\qquad + \mu^4\Bigl\{-\tfrac{2(m-2)(m-3)}{(m-1)^2}\,\Phi + \Omega_4\,\tfrac{(m-5)}{(m-1) H}\,\Omega_2 - \tfrac{1}{(m-1)^2}\,\tfrac{1}{H}\,\Omega_2^2\Bigr\}\Bigr]\\
            &= \tfrac{1}{9(m-1)^2 H^2}\bigl(4 W - 3 \Omega H\bigr),
\end{aligned}
\]
and consequently
\[
  \Lambda = \tfrac{1}{9 H^2}\bigl(4 W - 3 \Omega H\bigr),
\]
which agrees with the result before obtained, and thus the present method gives the complete solution of the problem.

\bigskip

% ============================================================
% PAPER 262
% ON THE EQUATION OF DIFFERENCES FOR AN EQUATION OF ANY ORDER,
% AND IN PARTICULAR FOR THE EQUATIONS OF THE ORDERS TWO, THREE,
% FOUR, AND FIVE.
% ============================================================

\newpage
\noindent\textbf{240}\hfill\textbf{[262}

\vspace{2cm}
\begin{center}
{\large\textbf{262.}}\\[1.5em]
{\large\textsc{On the Equation of Differences for an Equation of Any Order, and in particular for the Equations of the Orders Two, Three, Four, and Five.}}\\[1.5em]
\end{center}

\medskip
\noindent[\textit{From the Philosophical Transactions of the Royal Society of London}, vol.~\textsc{cl}. (for the year 1860), pp.~93--112. Received March 2,---Read March 29, 1860.]

\bigskip

\noindent The term, \textit{equation of differences}, denotes the equation for the squared differences of the roots of a given equation; the equation of differences afforded a means of determining the number of real roots, and also limits for the real roots, of a given numerical equation, and was upon this account long ago sought for by geometers. In the \textit{Philosophical Transactions} for 1763, Waring gives, but without demonstration or indication of the mode of obtaining it, the equation of differences for an equation of the fifth order wanting the second term: the result was probably obtained by the method of symmetric functions. This method is employed in the \textit{Meditationes Algebraic\ae} (1782), where the equation of differences is given for the equations of the third and fourth orders wanting the second terms; and in p.~85 the before-mentioned result for the equation of the fifth order wanting the second term, is reproduced. The formul\ae{} for obtaining by this method the equation of differences, are fully developed by Lagrange in the \textit{Trait\'e des Equations Num\'eriques} (1808); and he finds by means of them the equation of differences for the equations of the orders two and three, and for the equation of the fourth order wanting the second term; and in Note III. he gives, after Waring, the result for the equation of the fifth order wanting the second term. It occurred to me that the equation of differences could be most easily calculated by the following method. The coefficients of the equation of differences, \textit{qua} functions of the differences of the roots of the given equation, are leading

\newpage
\noindent\textbf{262]}\hfill\textsc{on the equation of differences for an equation of any order.}\hfill\textbf{241}

\medskip

\noindent coefficients of covariants, or (to use a shorter expression) they are ``Seminvariants\footnote{The term ``Seminvariant'' seems to me preferable to M.~Brioschi's term ``Peninvariant.''},'' that is, each of them is a function of the coefficients which is reduced to zero by one of the two operators which reduce an invariant to zero. In virtue of this property they can be calculated, when their values are known for the particular case in which one of the coefficients of the given equation is zero. To fix the ideas, let the given equation be $(*\Cdel v, 1)^n = 0$; then, when the last coefficient or constant term vanishes, the equation breaks up into $v = 0$ and into an equation of the degree $n - 1$, which I call the reduced equation; the equation of differences will break up into two equations, one of which is the equation of differences for the reduced equation, the other is the equation for the squares of the roots of the same reduced equation. This hardly requires a proof; let the roots of the given equation be $\alpha$, $\beta$, $\gamma$, $\delta$, \&c., those of the equation of differences are $(\alpha - \beta)^2$, $(\alpha - \gamma)^2$, $(\alpha - \delta)^2$, \&c., $(\beta - \gamma)^2$, $(\beta - \delta)^2$, $(\gamma - \delta)^2$, \&c.; but in putting the constant term equal to zero, we in effect put one of the roots, say $\alpha$, equal to zero; the roots of the equation of differences thus become $\beta^2$, $\gamma^2$, $\delta^2$, \&c., $(\beta - \gamma)^2$, $(\beta - \delta)^2$, $(\gamma - \delta)^2$, \&c. The equation for the squares of the roots can be found without the slightest difficulty; hence if the equation of differences for the reduced equation of the order $n - 1$ is known, we can, by combining it with the equation for the squares of the roots, form the equation of differences for the given equation with the constant term put equal to zero, and thence by the above-mentioned property of the Seminvariancy of the coefficients, find the equation of differences for the given equation. The present memoir shows the application of the process to equations of the orders two, three, four, and five: part of the calculation for the equation of the fifth order was kindly performed for me by the Rev.~R.~Harley. It is to be noticed that the best course is to apply the method in the first instance to the forms $(a, b, \ldots\Cdel v, 1)^n = 0$, without numerical coefficients (or, as they may be termed, the \textit{denumerate forms}), and to pass from the results so obtained to those which belong to the forms $(a, b, \ldots \mid v, 1)^n = 0$, or \textit{standard forms}. The equation of differences, for $(\alpha - \beta)^2$, \&c., the coefficients of which are seminvariants, naturally leads to the consideration of a more general equation having for its roots $(\alpha - \beta)^2 (x - \gamma y)\,(x - \delta y) \ldots$, \&c., the coefficients of which are covariants; and in fact, when, as for equations of the orders two, three, and four, all the covariants are known, such covariant equation can be at once formed from the equation of differences; for equations of the fifth order, however, where the covariants are not calculated beyond a certain degree, [they are now all calculated, see 141 and 143], only a few of the coefficients of the covariant equation can be thus at once formed. At the conclusion of the memoir, I show how the equation of differences for an equation of the order $n$ can be obtained by the elimination of a single quantity from two equations each of the order $n - 1$; and by applying to these two equations the simplification which I have made in B\'ezout's abridged method of elimination, I exhibit the equation of differences for the given equation of the order $n$, in a compendious form by means of a determinant; the first-mentioned method is, however, that which is best adapted for the actual development of the equation of differences for the equation of a given order.

\medskip

\hfill {\footnotesize\textsc{c. iv.}\hfill 31}

\newpage
\noindent\textbf{242}\hfill\textsc{on the equation of differences for an equation of any order.}\hfill\textbf{[262}

\medskip

\noindent The equations successively considered are
\[
\begin{aligned}
  (a, b, c \Cdel v, 1)^2 &= 0,\\
  (a, b, c, d \Cdel v, 1)^3 &= 0,\\
  (a, b, c, d, e \Cdel v, 1)^4 &= 0,\\
  (a, b, c, d, e, f \Cdel v, 1)^5 &= 0.
\end{aligned}
\]
The equation of differences for the quadric, and that for the squares of the roots, are considered to be known, and the other results are derived from them: it will be convenient to write down in the first instance the results for the quadric, the cubic, and the quartic equations, and then explain the process of obtaining them.

\medskip

\noindent For the quadric equation,
\[
  \text{Equation of differences is, } 0 = a^2(ac - b^2) + 4\,(\theta, 1)^1.
\]
That is, the equation of differences for $(a, b, c \Cdel v, 1)^2 = 0$ has the form
\[
  a^2\bigl(ac - b^2\bigr) + 4\theta = 0,
\]
and the equation for the squares of the roots is
\[
  (a^2,\; ac + 2b^2,\; c^2,\; ldots \Cdel \theta, 1)^2 = 0.
\]

\medskip
\noindent\textbf{Equation for the squares of the roots is, } $0 = $
\[
\begin{aligned}
  & a^2\quad (b^2 + 1)\\
  & +1\quad (ac + 2b^2)\quad (c^2 + 1)\quad (\theta, 1)^2.
\end{aligned}
\]

\medskip

\noindent For the cubic equation,
\[
  \text{Equation of differences is, } 0 = a^4(\theta, 1)^3,
\]
having coefficients (denumerate form)
\[
\begin{aligned}
  &+1,\\
  &a^2 c + 6,\quad a^2 c + 27,\quad ab^2 c + 96,\\
  &b^2 - 2,\quad b^2 + 1,\quad a^2 c^2 + 96,\\
  &\quad abcd + 18,\\
  &\quad ac^2 + 4,\\
  &\quad b^2 d + 4,\\
  &\quad b^2 c^2 - 1.
\end{aligned}
\]
The equation for the squares of the roots is
\[
\begin{aligned}
  & a^2 \times\\
  & +1,\quad ac + 2,\quad c^2 + 1,\quad d^2 + 1,\quad bd - 2,\quad (\theta, 1)^3.\\
  & b^2 - 1.
\end{aligned}
\]

\newpage
\noindent\textbf{262]}\hfill\textsc{on the equation of differences for an equation of any order.}\hfill\textbf{243}

\medskip

\noindent For the quartic equation,
\[
  \text{Equation of differences is, } 0 = a^4 \times \ldots\, (\theta, 1)^6,
\]
with denumerate-form coefficients listed in successive columns (each column is the coefficient of $\theta^k$ for $k = 0, 1, \ldots, 6$):
\medskip

\noindent\textit{Column } $a^4 \times$:
\[
  +1,\quad ac + 8,\quad b^2 - 3.
\]

\noindent\textit{Column } $a^4 \theta \times$:
\[
  a^2 e + 8,\quad a^2 ce + 16,\quad a^2 b d - 2,\quad a^2 c^2 + 22,\quad a c^2 - 16,\quad b^4 + 3.
\]

\noindent\textit{Column } $a^4 \theta^2 \times$:
\[
\begin{aligned}
  & a^2 ce + 16,\quad a^2 c^2 + 26,\\
  & a^2 b^2 e - 6,\quad a^2 b c d - 30,\quad a^2 c^2 - 16,\\
  & a^2 b^3 d + 288,\quad a^2 b^2 c^2 - 24,\quad a b^3 c + 8,\quad b^6 - 1.
\end{aligned}
\]

\noindent\textit{Column } $a^4 \theta^3 \times$:
\[
\begin{aligned}
  & a^2 b c^2 e + 48,\quad a^2 b^2 d e + 24,\quad a^2 c d^2 + 32,\quad a^2 b^2 c^2 + 17,\\
  & a^2 b^3 d e + 54,\quad a^2 b c d^2 + 32,\quad a b^3 c d^2 + 25,\quad a b^2 e + 38,\quad a b^3 d^2 + 6,\quad a b^2 d^2 + 26.
\end{aligned}
\]

\noindent\textit{Column } $a^4 \theta^4 \times$:
\[
\begin{aligned}
  & a^2 e^2 + 256,\quad a^2 b^2 c d - 192,\\
  & a^2 b^2 c^2 + 216,\quad a^2 b c^2 d e - 144,\quad a^2 b^2 d^2 - 128,\\
  & a^2 b c d e - 54,\quad a^2 b d^3 - 120,\quad a^2 d^4 - 27,\\
  & a b^3 c d - 144,\quad a b^3 d^2 - 6,\quad a b^2 c d - 80,\quad a b c d^3 + 18,\\
  & a c^4 d^2 + 16,\quad a c^3 d^2 - 4,\quad b^4 d^2 + 49,\quad b^4 c^2 - 27,\\
  & b^3 c d^2 + 18,\quad b^3 d^4 - 4,\quad b^2 c^4 - 1,\quad b^2 c d^2 + 1,\quad b^2 c^3 e - 4.
\end{aligned}
\]

\noindent (The full table is laid out in six columns on the printed page; the above lists the entries column by column, in the order of the printed table.)

\medskip

\noindent\textbf{Equation for squares of the roots is, } $0 = $
\[
\begin{aligned}
  & a^2 \times\\
  & +1,\quad ac + 2,\quad ae + 2,\quad ce + 2,\quad e^2 + 1,\\
  & b^2 - 1,\quad c^2 + 1,\quad bd - 2,\quad d^2 - 1,\quad (\theta, 1)^4.
\end{aligned}
\]

\medskip

\noindent The multiplication of the equation of differences and the equation of the squares of the roots of the quadric equation, gives the equation, $0 = (\theta, 1)^3$,
\[
\begin{aligned}
  & a^4 \times\\
  & +1,\quad ac + 6,\quad a^2 c^2 + 9,\quad ac^3 + 4,\\
  & b^2 - 4,\quad ab^2 c + 16,\quad b^2 c^2 - 1,\quad b^4 + 4,\\
\end{aligned}
\]
where all the coefficients except the last are reduced to zero by the operator $(\theta, 1)^3$,
\[
  3a\partial_b + 2b\partial_c + c\partial_d,
\]

\medskip

\hfill {\footnotesize 31--2}

\newpage
\noindent\textbf{244}\hfill\textsc{on the equation of differences for an equation of any order.}\hfill\textbf{[262}

\medskip

\noindent and they are consequently (without any alteration) coefficients of the equation of differences of the cubic equation: the last coefficient is not reduced to zero by the operator, and requires therefore to be completed by the adjunction of the terms in $d$ (the series, here and in every other case, is of course a finite one, the number of terms might easily be calculated \textit{\`a priori}). Let the value be $L_0 + L_1 d + L_2 d^2 + \&\text{c.}$, we have $L_0 = +4ac^3 - b^2 c^2$; and putting for shortness $\nabla = 3a\partial_b + 2b\partial_c$, the operator which reduces this to zero is $\nabla + c\partial_d$; we ought therefore to have
\[
\begin{aligned}
  0 &= \nabla L_0 + d(\nabla L_1 + \ldots) + d^2(\nabla L_2 + \ldots)\\
    &\quad + cL_1 + 2cL_2 + 3cL_3,
\end{aligned}
\]
and consequently
\[
  L_1 = -\tfrac{1}{c}\nabla L_0,\quad L_2 = -\tfrac{1}{2c}\nabla L_1,\quad L_3 = -\tfrac{1}{3c}\nabla L_2,\quad \&\text{c.},
\]
giving
\[
  L_1 = -\tfrac{1}{c}\bigl\{(-6 + 24c^2)\,18 abc^2 - 4 b^3 c\bigr\} = -18abc + b^3,
\]
\[
  L_2 = -\tfrac{1}{2c}\bigl\{-54a^2 c + (36 - 36c^2)\,0\,ab^2\bigr\} = +27 a^2,
\]
\[
  L_3 = 0,\quad \&\text{c.},
\]
and consequently for the last coefficient the value above written down; it will be presently seen how in more complicated cases the calculations should be arranged.

\medskip

\noindent Again, multiplying together the equation of differences and the equation for the squares of the roots of the cubic equation, we obtain an equation which it is not necessary to write down, as it can be at once formed by putting $e = 0$ in the equation of differences for the quartic equation. And from the equation so obtained, by the adjunction of the terms in $e$, we find the equation of differences for the quartic equation, viz. each coefficient is of the form $L_0 + L_1 e + L_2 e^2 + \&\text{c.}$, where $L_0$ is known, and such coefficient is reduced to zero by the operator
\[
  4a\partial_b + 3b\partial_c + 2c\partial_d + d\partial_e,
\]
or putting for shortness $\nabla = 4a\partial_b + 3b\partial_c + 2c\partial_d$, the operator $\nabla + d\partial_e$. We have therefore
\[
  L_1 = -\tfrac{1}{d}\nabla L_0,\quad L_2 = -\tfrac{1}{2d}\nabla L_1,\quad L_3 = -\tfrac{1}{3d}\nabla L_2,\quad \&\text{c.}
\]
It is to be observed that the last coefficient of the equation of differences is the discriminant, and that the above method of calculating the coefficients of the equation of differences, as applied to the last coefficient, is nothing else than the method of calculating the discriminant given in my Fourth Memoir on Quantics, [155].

\newpage
\noindent\textbf{262]}\hfill\textsc{on the equation of differences for an equation of any order.}\hfill\textbf{245}

\medskip

\noindent The multiplication of the equation of differences, and the equation for the squares of the roots of the quartic equation, gives, in like manner, the equation of differences for the quintic equation, except as to the terms involving $f$; and these are obtained as above, viz. each coefficient is of the form $L_0 + L_1 f + L_2 f^2 + \&\text{c.}$, where $L_0$ is known; such coefficient is reduced to zero by the operator
\[
  5a\partial_b + 4b\partial_c + 3c\partial_d + 2d\partial_e + e\partial_f,
\]
or putting for shortness $\nabla' = 5a\partial_b + 4b\partial_c + 3c\partial_d + 2d\partial_e$, the operator $\nabla' + e\partial_f$. We therefore have
\[
  L_1 = -\tfrac{1}{e}\nabla' L_0,\quad L_2 = -\tfrac{1}{2e}\nabla' L_1,\quad L_3 = -\tfrac{1}{3e}\nabla' L_2,\quad \&\text{c.},
\]
which give $L_1$, $L_2$, \&c. by means of $L_0$. For the calculation of $\nabla L$ (where $L$ is any one of the coefficients $L_0$, $L_1$, \&c.), it is proper to separate the terms involving the different powers of $e$, and write $L = M_0 + M_1 e + M_2 e^2 + \&\text{c.}$, $\nabla = \nabla'' + 2d\partial_e$, where $\nabla'' = 5a\partial_b + 4b\partial_c + 3c\partial_d$. We have then
\[
  \nabla L = \nabla'' M_0 + e(\nabla'' M_1 + 2 d M_2) + \&\text{c.}\\
            + 2 d M_1,
\]
or, what is the same thing,
\[
  \frac{1}{1}\,\nabla L_n + \frac{4 d M_1}{\ldots} + \ldots,
\]
and as an equation which should be satisfied identically, and which would therefore serve as a verification,
\[
  \nabla'' M_n + 2 d M_1 = 0;
\]
but, since a verification was obtained by other means, the equations of this kind were not used for the purpose. It may be interesting to give the actual calculation of one of the coefficients, say of coefficient $\theta^4$ (which, with coefficient $\theta^5$, was calculated by Mr~Harley).

\newpage
\noindent\textbf{246}\hfill\textsc{on the equation of differences for an equation of any order.}\hfill\textbf{[262}

\medskip

\noindent\textit{Calculation of coefficient $\theta^4$ in equation of differences for the quintic equation}
\[
  (a, b, c, d, e, f \Cdel v, 1)^5 = 0.
\]

\medskip

\noindent\textit{Calculation of $L_0$.}

\medskip

\noindent [The original page presents a large computational table giving the coefficients of $L_0$ as a sum
\[
  L_0 = A + Be + Ce^2 + De^3 + Ee^4,
\]
with each of $-A$, $-B$, $-C$, $-D$, $-E$ expanded as a sum of monomial terms in $a$, $b$, $c$, $d$, $e$, $f$ with integer numerical coefficients. The principal entries, transcribed column by column from the original tabular layout, are:]

\medskip

\noindent$-E = $
\[
\begin{aligned}
  & a^3\,c d^2 \cdot (-960),\quad a^3 b d^2 \cdot (+840),\quad a^4 \cdot (-400),\\
  & a^2 d^3 \cdot (+270),\\
  & a^2 b c d \cdot (+240),\quad a^2 b^2 d^2 \cdot (-275),\\
  & a^2 b c^2 d \cdot (+390),\quad a^2 c^3 \cdot (+360),\quad a b^3 c \cdot (-400),\\
  & a^2 b^2 d^2 \cdot (-140),\quad a b^4 \cdot (+48),\\
  & a b^2 c d^2 \cdot (+14),\quad a^2 c^3 \cdot (+15),\\
  & \quad a b^3 c d \cdot (+178),\\
  & a b^2 c^2 d^2 \cdot (-246),\\
  & a b c^4 d \cdot (+164),\quad a b^2 c^3 \cdot (-112),\\
  & a c^5 \cdot (-24),\\
  & b^3 d^3 \cdot (-48),\\
  & b^2 c d^3 \cdot (+54),\quad b^4 c^2 \cdot (+25),\\
  & b^2 c^2 d \cdot (-38),\\
  & b^2 c^4 \cdot (+6).
\end{aligned}
\]

\medskip

\noindent (The full layout of the original table involves further columns headed $-D$, $-C$, $-B$, $-A$, each enumerating numerous monomials in $a$, $b$, $c$, $d$, $e$, $f$. Owing to the very large number of entries (several hundreds of monomials in five columns) and the dense matrix-style typesetting in the original, the table is reproduced here only in summary form; the reader is referred to the original page for the complete listing.)

\newpage
\noindent\textbf{262]}\hfill\textsc{on the equation of differences for an equation of any order.}\hfill\textbf{247}

\medskip

\noindent\textit{Calculation of $L_1$.}

\medskip

\noindent Terms of $L_1 f$ not involving $e$:
\[
  -f(\nabla'' B + 4 d C), \quad\text{viz.}\quad\text{(table of coefficients)}.
\]
[The original page gives a large rectangular table tabulating, for each monomial in $a$, $b$, $c$, $d$ multiplied by $f$, the contributions $5a\partial_b B$, $4b\partial_c B$, $3c\partial_d B$, $4 d C$, and the resulting coefficient of the monomial in $L_1$. Principal monomials listed include $a^2 c d^2 f$, $a^2 b^2 d^2 f$, $a^2 b c^2 d^2 f$, $a b^3 c d^2 f$, $a b c^4 d f$, $a c^5 f$, $b^4 d^2 f$, $b^2 c d^3 f$, $b^2 c d f$, $b^2 c^2 d f$ with respective totals reported in the rightmost column ($-420$, $+600$, $+5070$, $+24$, $+40$, $+130$, $+32$, $+24$, $+6$, etc.).]

\medskip

\noindent Terms of $L_1 f$ involving $e$:
\[
  -e f(\nabla'' C + 6 d D), \quad\text{viz.}\quad\text{(table of coefficients)}.
\]
[Likewise tabulated; principal monomials include $a^2 c d^2 e f$, $a^2 b^2 d e f$, $a^2 b c d e f$, $a^2 b c^2 e f$, $a b^3 d e f$, $a b^2 c^2 e f$, $b^4 c e f$ with totals $-1550$, $+1740$, $-200$, $+290$, $+310$, $+288$, $+56$, etc.]

\medskip

\noindent Terms of $L_1 f$ involving $e^2$:
\[
  -e^2 f(\nabla'' D + 8 d E), \quad\text{viz.}\quad\text{(table of coefficients)}.
\]
[Principal monomials include $a^2 d^3 f$, $a^2 b c d f$ with totals $+1000$, $+1400$, etc.]

\newpage
\noindent\textbf{248}\hfill\textsc{on the equation of differences for an equation of any order.}\hfill\textbf{[262}

\medskip

\noindent\textit{Calculation of $L_2$.}

\medskip

\noindent $L_2 = A + B e + C e^2,\quad\text{suppose, where}$
\[
\begin{aligned}
  -A &:\quad a^2 b d^2 \cdot (+775),\quad a^4 d \cdot (-500),\quad a^3 b c \cdot (-700),\quad a^3 b^2 \cdot (+320),\\
       &\quad a^2 b^2 c d \cdot (-870),\\
       &\quad a^2 b c^2 \cdot (-145),\\
       &\quad a b^4 d \cdot (+100),\\
       &\quad a b^3 c \cdot (+155),\\
       &\quad b^5 \cdot (-28),\\
  -B &:\quad a^2 b c d f^2 \cdot (\ldots).
\end{aligned}
\]

\medskip

\noindent Terms of $L_2 f^2$ not involving $e$ are given by
\[
  -f^2(\nabla'' B + 4 d^2 C),
\]
tabulated as above. Principal contributions yield (for example)
\[
  a^3 b c d f^2 \cdot (+1875),\quad a^2 c d^2 f^2 \cdot (+200),\quad a^2 b^2 d f^2 \cdot (-2025),\quad a^2 c^3 f^2 \cdot (+1750),\quad a b^4 d f^2 \cdot (-112),\quad a b^3 c d f^2 \cdot (+840).
\]

\medskip

\noindent Terms of $L_2 f^2$ involving $e$ are given by
\[
  -e f^2(\nabla'' C + 6 d D),
\]
with principal contributions $a^2 c e f^2 \cdot (-5000)$, $a^2 d e f^2 \cdot (+2000)$, etc.

\newpage
\noindent\textbf{262]}\hfill\textsc{on the equation of differences for an equation of any order.}\hfill\textbf{249}

\medskip

\noindent\textit{Calculation of $L_3$\quad(gives $L_3 = 0$).}

\medskip

\noindent $L_3 = A + B e$, where
\[
  -A:\quad a^3 c\cdot(+1666\tfrac{2}{3}),\quad a^2 b^2\cdot(-666\tfrac{2}{3});
\]
and
\[
  -B:\quad\text{is determined from } -\tfrac{1}{3}\,\nabla'' B \text{ via the table; the resulting coefficient is } 0.
\]

\medskip

\noindent Terms of $L_3 f^3$ are given by
\[
  -\tfrac{1}{3}\,f^3\,\nabla'' B \quad\text{viz.}\quad a^3 b f^3\cdot(-6666\tfrac{2}{3}) + (+6666\tfrac{2}{3}) = 0.
\]

\medskip

\noindent and the required coefficient of $\theta^4$ is
\[
  L_0 + L_1 f + L_2 f^2.
\]

\medskip

\noindent All the coefficients were calculated in this manner, except the last coefficient, which was deduced from the known value of the discriminant for the standard form. And we have thus the complete expression of the equation of differences for the general quintic equation under the denumerate form $(a, b, c, d, e, f \Cdel v, 1)^5 = 0$, viz.

\medskip

\hfill {\footnotesize\textsc{c. iv.}\hfill 32}

\newpage
\noindent\textbf{250--251}\hfill[Two-page tabular display]\hfill\textbf{[262}

\medskip

\noindent\textit{Equation of differences for $(a, b, c, d, e, f \Cdel v, 1)^5 = 0$ is}
\medskip

\noindent [The original printed pages 250--251 present, across a double-page spread, the complete equation of differences for the general quintic in denumerate form. The result is a polynomial in $\theta$ of degree~6, in which each coefficient is itself a large polynomial in $a$, $b$, $c$, $d$, $e$, $f$ (degree~$2k$ in the entries for the coefficient of $\theta^{6-k}$, by virtue of the seminvariancy property). The seven columns are headed by the successive powers $\theta^0$, $\theta^1$, $\theta^2$, $\theta^3$, $\theta^4$, $\theta^5$, $\theta^6$ (with the appropriate multiplier $a^{2k}$ prefixed). Each column is a long list of monomials with their integer coefficients; principal entries include (in successive columns):]
\medskip

\noindent\textbf{Coefficient of $\theta^0$:}\quad $a^4\bigl(+1,\; ac+10,\; a^2 d - 4,\; b^2 - 4,\; ldots\bigr)$, with further entries $a^2 e + 30$, $a c^2 + 50$, $a^2 d c + 25$, $a^2 b e - 20$, $a b^2 c + 30$, $a^3 b d + 81$, $b^6 + \ldots$, etc.

\medskip

\noindent\textbf{Coefficient of $\theta^1$:}\quad $a^4\,\theta\bigl(a^2 f + \ldots, ab e f + 100, a^2 b d e - 110, \ldots\bigr)$.

\medskip

\noindent\textbf{Coefficient of $\theta^2$:}\quad $a^4\,\theta^2 \bigl( a^3 b e f - 250, a^2 d^2 e - 360, a^3 e^2 + 950, a^4 b d f + 625, \ldots \bigr)$ together with very many further monomials.

\medskip

\noindent\textbf{Coefficient of $\theta^3$:}\quad $a^4\,\theta^3\,(\ldots)$.

\medskip

\noindent\textbf{Coefficient of $\theta^4$:}\quad $a^4\,\theta^4\,(\ldots)$, the result of the calculation given in the preceding section.

\medskip

\noindent\textbf{Coefficient of $\theta^5$:}\quad $a^4\,\theta^5\,(\ldots)$.

\medskip

\noindent\textbf{Coefficient of $\theta^6$:}\quad $a^4\,\theta^6\,(\ldots)$, equal to the discriminant of the quintic up to the standard normalisation.

\medskip

\noindent (Owing to the very large extent of this table---more than four hundred monomials in six variables---the entries are not reproduced individually in the present transcription. The reader is referred to the original double-page spread of the printed memoir for the full list of monomials and coefficients.)

\newpage
\noindent\textbf{252}\hfill\textsc{on the equation of differences for an equation of any order.}\hfill\textbf{[262}

\medskip

\noindent It may be remarked, that if $\omega$ is an imaginary cube root of unity, then the roots of the equation $(1, 1, 1, 1, 1, 1 \Cdel v, 1)^5 = 0$ are $-1$, $\omega$, $\omega^2$, $-\omega$, $-\omega^2$; the differences of the roots are
\[
\begin{aligned}
  & -1 - \omega,\ -1 - \omega^2,\ -1 + \omega,\ -1 + \omega^2,\ \omega - \omega^2,\ 2\omega,\ \omega + \omega^2,\ \omega^2 + \omega,\ 2\omega^2,\ -\omega + \omega^2,\\
  =\,& \omega^2,\ \omega,\ -1 + \omega,\ -1 + \omega^2,\ \omega - \omega^2,\ 2\omega,\ -1,\ -1,\ 2\omega^2,\ -\omega + \omega^2,
\end{aligned}
\]
and the squares of the differences are
\[
  \omega,\ \omega^2,\ -3\omega,\ -3\omega^2,\ -3,\ 4\omega^2,\ 1,\ 1,\ 4\omega,\ -3,
\]
from which the equation of differences is found to be
\[
  (\theta^2 + \theta + 1)(\theta^2 - 3\theta + 9)(\theta^2 + 3\theta + 9)(\theta^2 - 2\theta + 1)(\theta^2 + 6\theta + 9) = 0;
\]
or multiplying out, it is
\[
  (1,\ 6,\ 21,\ 46,\ 108,\ 546,\ 493,\ -1410,\ -567,\ -540,\ +12960 \Cdel \theta,\ 1)^{10} = 0;
\]
which is what the preceding expression of the equation of differences becomes upon writing therein $a = b = c = d = e = f = 1$. Moreover, upon passing (as will presently be done) to the standard form, and then writing $a = b = c = d = e = f = 1$, all the coefficients (except the first coefficient, which is equal to unity) should become equal to zero; these two tests afford a complete verification of the result.

\medskip

\noindent The following corrections have to be made in Waring's result, as given by himself and Lagrange (Waring, \textit{Phil. Trans.} 1762).

\medskip

\noindent Waring, \textit{Meditationes Algebraic\ae}, p.~85:
\[
  \text{for } +169\,q^2 \text{ read } +196\,q^2 \quad \text{(in coefficient } w^2\text{)}.
\]

\noindent Lagrange, \textit{Equations Num\'eriques}, p.~108:
\[
\begin{aligned}
  \text{for } &+1200\,GE          \text{ read } +200\,GE   & &\text{(in } d\text{)},\\
  \text{for } &-169\,B^2 D        \text{ read } -196\,B^2 D & &\text{(in } e\text{)},\\
  \text{for } &-25\,B^4           \text{ read } +25\,B^4    & &\text{(in } f\text{)},\\
  \text{for } &+27\,G^2 B^2       \text{ read } -27\,G^2 D^2 & &\text{(in } k\text{)}.
\end{aligned}
\]

\medskip

\noindent It may be noticed, that if in the coefficients of the several powers of $\theta$ (as they are written down in the columns, without regarding the power of $a$ which multiplies the entire column), we attend only to the terms independent of $a$, we have the series
\[
\begin{aligned}
  &1,\ b^2 - 4,\ b^4 + 6,\ b^6 - 4,\ b^8 + 1,\ b^5 d + 8,\ b^6 f + 8,\ \&\text{c.,}\\
  &\hspace{12em} b^4 c e + 22,\\
  &\hspace{8em} b^5 c^2 - 3,\quad b^4 d^2 - 2,\\
  &\hspace{12em} b^4 c^2 d - 16,\\
  &\hspace{12em} b^4 c^4 + 3,
\end{aligned}
\]
the law of the first terms of which, up to the term $+ b^8$, is obvious; but the term $+ b^8$, which is the last term of this initial series, is also the first term of a terminal series, the terms of which are deduced from the coefficients in the equation of differences for the quartic equation $(a, b, c, d, e \Cdel v, 1)^4 = 0$, viz. these coefficients are

\newpage
\noindent\textbf{262]}\hfill\textsc{on the equation of differences for an equation of any order.}\hfill\textbf{253}

\medskip
\[
\begin{aligned}
  & a^4 \times\\
  & +1,\quad ac + 8,\quad a^2 e + 8,\\
  & b^2 - 3,\quad a^2 b d - 2,\\
  & \quad a^2 c^2 + 22,\\
  & \quad abc - 16,\\
  & \quad b^4 + 3;
\end{aligned}
\]
and by writing $b$, $c$, $d$, $e$, $f$ in the place of $a$, $b$, $c$, $d$, $e$ respectively, and multiplying by $b^4$, we have the above-mentioned series,
\[
  b^4 + 1,\ b^5 d + 8,\ \&\text{c.,}\quad b^5 c^2 - 3.
\]
It is easy to see, \textit{\`a priori}, in the case of an equation of any order, that this property holds good.

\medskip

\noindent Passing now to the standard forms:

\medskip

\noindent For the quadric $(a, b, c \Cdel v, 1)^2 = 0$, the equation of differences is, $0 = $
\[
\begin{aligned}
  & a^2 \times\\
  & +1\quad \text{and}\quad 4\times(ac + 1)\,(\theta,\,1)^1.
\end{aligned}
\]

\medskip

\noindent For the cubic equation $(a, b, c, d \Cdel v, 1)^3 = 0$, the equation of differences is, $0 = $
\[
\begin{aligned}
  & a^4 \times\\
  & +1,\quad 18\,a^2 \times (ac + 1,\; b^2 - 1),\\
  & 81 \times (a^2 c^2 + 1,\; b^2 a b d + 21),\\
  & 27 \times (a^2 d^2 + 1,\; abcd - 6,\; ac^3 + 4,\; b^3 d + 4,\; b^2 c^2 - 3)\quad (\theta, 1)^3.
\end{aligned}
\]

\medskip

\noindent For the quartic equation $(a, b, c, d, e \Cdel v, 1)^4 = 0$, the equation of differences is, $0 = $

\medskip

\noindent (The standard-form expression has the same monomial structure as the denumerate form above, but with the entries multiplied by the appropriate binomial-coefficient powers --- specifically $48\,a^4 \times$, $8\,a^4 \times$, $32 \times$, $16 \times$, $1152 \times$, $256 \times$ as overall numerical factors for the successive columns. The full standard-form quartic expression is set out as a six-column table on page~253 of the original, with principal entries including:
\[
\begin{aligned}
  & a^4 \times +1,\quad ac + 1,\quad a^2 e + 1,\quad a^2 ce + 3,\quad a^2 d^2 + 13,\\
  & b^2 - 1,\quad a^2 b d - 4,\quad a^2 b^2 e + 56,\quad a^3 b^2 c d e - 7,\quad a^3 e^2 - 1,\\
  & a^2 c^2 + 99,\quad a^2 c d^2 + 288,\quad a^2 d^2 e + 3,\quad a^2 c^2 e^2 + 1,\\
  & a^2 b^2 c e - 192,\quad a^2 b^2 d^2 - 400,\quad a^2 b c d e - 10,\quad a^2 c d^2 e + 54,\\
  & a^2 b c^2 - 192,\quad a^2 b c^2 d - 1944,\quad a^2 b d^3 + 1269,\quad a^2 c^3 d^2 + 54,\\
  & b^4 + 96,\quad a^2 c^4 + 1377,\quad a b^3 d^2 - 6,\quad a b^2 c d^2 + 27,\\
  & a^2 b^3 d + 64,\quad a b^3 c d - 144,\quad a b^3 d^2 e - 6,\quad a b^2 c^3 e + 54,\\
  & a^2 b^4 c - 432,\quad a b^2 c d - 80,\quad a b^2 c^2 d - 78,\quad a b c d^3 + 108,\\
  & a b^3 c + 384,\quad a c^4 + 1377,\quad a b^2 c^3 - 78,\quad a c^3 e + 81,\\
  & b^6 - 128,\quad a b^4 e + 96,\quad a c^4 + 27,\quad a c^2 d^2 - 54,\\
  & \quad a b^3 c d + 3648,\quad a b^2 c d^3 + 4,\quad b^4 c e + 108,\\
  & \quad a b^4 c - 2592,\quad a b c d^2 - 180,\quad b^4 d^2 + 27,\\
  & \quad b^4 d - 1536,\quad a b^3 c^2 + 108,\quad b^3 c d^2 + 18,\\
  & \quad b^6 + 1152,\quad a c^4 d + 81,\quad b^2 c d^4 + 36;\\[4pt]
  & \quad\text{etc.},\quad (\theta, 1)^4.
\end{aligned}
\]
Owing to the very dense layout of the original printed table, only the principal entries are reproduced above; the reader is referred to the original page 253 of the memoir for the complete six-column table of monomials and coefficients.)

\medskip

\noindent [The paper continues on the following pages 254 onwards with the standard-form expression for the quintic and the alternative determinant-elimination method; these pages fall outside the present extract.]

\end{document}
